\clearpage\pagestyle{otMain}\fontsize{11}{15}\selectfont\corrigesChapitre{16}{Espace, solides et repérage sur la sphère}
\begin{blocLivre}{cor-16-1}\begin{corrige}{1}{Un pavé}\demarche{Multiplier les trois dimensions perpendiculaires.}
$V=8\times5\times3=\res{120\text{ cm}^3}$. Les trois dimensions sont exprimées dans la même unité.
\label{fig-cor-16-1}\figureGeo[corrige]{espace-01}
\end{corrige}\end{blocLivre}
\begin{blocLivre}{cor-16-2}\begin{corrige}{2}{Un cylindre}\demarche{Calculer l’aire du disque de base.}
$V=\pi\times3^2\times10=\res{90\pi\text{ cm}^3}\approx\res{283\text{ cm}^3}$. Le rayon est déjà fourni, il ne faut pas le diviser par deux.
\label{fig-cor-16-2}\figureGeo[corrige]{espace-02}
\end{corrige}\end{blocLivre}
\begin{blocLivre}{cor-16-3}\begin{corrige}{3}{Une pyramide}\demarche{Utiliser la hauteur perpendiculaire.}
L’aire de base vaut 36 cm$^2$. Le volume vaut $36\times9/3=\res{108\text{ cm}^3}$. La division par trois distingue la pyramide du prisme de même base et hauteur.
\label{fig-cor-16-3}\figureGeo[corrige]{espace-03}
\end{corrige}\end{blocLivre}
\begin{blocLivre}{cor-16-4}\begin{corrige}{4}{Un cône}\demarche{Passer du diamètre au rayon.}
Le rayon est 4 cm. $V=\pi\times4^2\times9/3=\res{48\pi\text{ cm}^3}$. Employer 8 comme rayon multiplierait le résultat par quatre.
\label{fig-cor-16-4}\figureGeo[corrige]{espace-04}
\end{corrige}\end{blocLivre}
\begin{blocLivre}{cor-16-5}\begin{corrige}{5}{Une boule}\demarche{Ne pas confondre aire et volume.}
$V=4\pi\times27/3=\res{36\pi\text{ cm}^3}$. L’aire vaut $4\pi\times9=\res{36\pi\text{ cm}^2}$. Les valeurs numériques des coefficients coïncident ici, mais les grandeurs et unités diffèrent.
\label{fig-cor-16-5}\figureGeo[corrige]{espace-05}
\end{corrige}\end{blocLivre}
\begin{blocLivre}{cor-16-6}\begin{corrige}{6}{Lire un patron de cube}\demarche{Simuler le pliage sans confondre proximité sur la feuille et sur le cube.}
La bande horizontale B-C-D-E s’enroule : B est opposée à D, C à E. A et F ferment les extrémités et sont opposées. L’aire totale vaut $6\times2^2=\res{24\text{ cm}^2}$.
\label{fig-cor-16-6}\figureGeo[corrige]{espace-patron}
\end{corrige}\end{blocLivre}
\begin{blocLivre}{cor-16-7}\begin{corrige}{7}{Trois vues}\demarche{Associer chaque vue aux dimensions visibles.}
La face montre $\res{8\times3}$ cm, le dessus $\res{8\times5}$ cm et le côté $\res{5\times3}$ cm. La profondeur n’apparaît pas dans la vue de face.
\label{fig-cor-16-7}\figureGeo[corrige]{espace-vues}
\end{corrige}\end{blocLivre}
\begin{blocLivre}{cor-16-8}\begin{corrige}{8}{Une section de pavé}\demarche{Reconnaître une section parallèle à une face.}
La section est un rectangle de dimensions 8 cm sur 5 cm. Son aire vaut $\res{40\text{ cm}^2}$. La hauteur de coupe ne modifie pas cette section parallèle dans un pavé.
\label{fig-cor-16-8}\figureGeo[corrige]{espace-section-pave}
\end{corrige}\end{blocLivre}
\begin{blocLivre}{cor-16-9}\begin{corrige}{9}{Des coordonnées dans un pavé}\demarche{Identifier le rôle des trois coordonnées.}
Le sommet opposé a pour coordonnées $\res{(6;4;3)}$. La projection verticale garde les deux premières coordonnées et annule la hauteur : $\res{(2;3;0)}$. P est bien à l’intérieur des trois limites.
\label{fig-cor-16-9}\figureGeo[corrige]{espace-repere}
\end{corrige}\end{blocLivre}
\begin{blocLivre}{cor-16-10}\begin{corrige}{10}{Latitude et longitude}\demarche{Distinguer latitude, longitude et distance.}
Ils ont la même latitude : même parallèle. Leurs longitudes diffèrent : méridiens distincts. Ils sont de part et d’autre de l’origine des longitudes, donc l’écart vaut $20+50=\res{70^\circ}$. Cet angle n’est pas une distance en kilomètres.
\label{fig-cor-16-10}\figureGeo[corrige]{espace-sphere}
\end{corrige}\end{blocLivre}
\begin{blocLivre}{cor-16-11}\begin{corrige}{11}{Un réservoir assemblé}\demarche{Additionner les volumes puis convertir.}
Le cylindre vaut $2\pi$ m$^3$, la demi-boule $2\pi/3$ m$^3$. Le total est $8\pi/3$ m$^3$, soit $8000\pi/3$ L, environ $\res{8378\text{ L}}$. Les deux parties ne se recouvrent qu’à leur frontière.
\label{fig-cor-16-11}\figureGeo[corrige]{espace-reservoir}
\end{corrige}\end{blocLivre}
\begin{blocLivre}{cor-16-12}\begin{corrige}{12}{Une boule creusée}\demarche{Soustraire les volumes du plein et du creux.}
On soustrait les volumes : $V=4\pi(5^3-2^3)/3=4\pi\times117/3=\res{156\pi\text{ cm}^3}$. Il faut soustraire des cubes, pas calculer le volume d’un rayon $5-2$.
\label{fig-cor-16-12}\figureGeo[corrige]{espace-cavite}
\end{corrige}\end{blocLivre}
\begin{blocLivre}{cor-16-13}\begin{corrige}{13}{Une section de sphère}\demarche{Passer par le triangle rectangle de la section.}
Pour M sur le cercle, OHM est rectangle en H. $HM^2=OM^2-OH^2=25-9=16$, donc $\res{HM=4\text{ cm}}$. Le disque a pour aire $\res{16\pi\text{ cm}^2}$.
\label{fig-cor-16-13}\figureGeo[corrige]{espace-section-sphere}
\end{corrige}\end{blocLivre}
\begin{blocLivre}{cor-16-14}\begin{corrige}{14}{Un cône réduit}\demarche{Utiliser le cube du rapport puis soustraire.}
Le grand volume vaut $\pi\times36\times12/3=144\pi$ cm$^3$. Le petit vaut $(1/2)^3\times144\pi=\res{18\pi\text{ cm}^3}$. Le tronc vaut $144\pi-18\pi=\res{126\pi\text{ cm}^3}$.
\label{fig-cor-16-14}\figureGeo[corrige]{espace-cone-coupe}
\end{corrige}\end{blocLivre}
\begin{blocLivre}{cor-16-15}\begin{corrige}{15}{La diagonale de l’espace}\demarche{Construire d’abord une diagonale de face.}
La diagonale de base vaut $\sqrt{3^2+4^2}=5$ cm. Avec la hauteur de 12 cm, un autre triangle rectangle donne $\sqrt{5^2+12^2}=\res{13\text{ cm}}$. Deux applications successives de Pythagore suffisent.
\label{fig-cor-16-15}\figureGeo[corrige]{espace-diagonale}
\end{corrige}\end{blocLivre}
